% ═══════════════════════════════════════════════════════════════
% MT-06d-ML: Minitest MUSTERLÖSUNG - Mi ciudad (Wiederholung)
% Spanisch Klasse 7 | Sequenz 6: Mi ciudad (Abschluss)
% ═══════════════════════════════════════════════════════════════

\documentclass[11pt, a4paper]{article}

\usepackage[utf8]{inputenc}
\usepackage[T1]{fontenc}
\usepackage[ngerman]{babel}
\usepackage{helvet}
\renewcommand{\familydefault}{\sfdefault}

\usepackage[margin=2cm]{geometry}
\usepackage{enumitem}
\usepackage{tabularx}
\usepackage{booktabs}

\usepackage{xcolor}
\usepackage{tcolorbox}
\tcbuselibrary{skins}

\usepackage{fancyhdr}
\usepackage{lastpage}
\usepackage{needspace}

% FARBEN
\definecolor{spanisch-color}{HTML}{c41e3a}
\definecolor{grau-color}{HTML}{888888}
\definecolor{hellgrau-color}{HTML}{f5f5f5}
\definecolor{loesung-color}{HTML}{c62828}

\colorlet{spanisch}{spanisch-color}
\colorlet{grau}{grau-color}
\colorlet{hellgrau}{hellgrau-color}
\colorlet{loesung}{loesung-color}

\newcommand{\fachfarbe}{spanisch}

\newcommand{\thema}{Minitest: Mi ciudad -- Wiederholung}
\newcommand{\fach}{Spanisch}
\newcommand{\klasse}{7}
\newcommand{\maxpunkte}{30}
\newcommand{\zeit}{20 Minuten}
\newcommand{\datum}{\today}

\pagestyle{fancy}
\fancyhf{}
\renewcommand{\headrulewidth}{0.4pt}
\renewcommand{\footrulewidth}{0.4pt}

\lhead{\fach{} -- Klasse \klasse}
\chead{\textcolor{loesung}{\textbf{MT-06d MUSTERLÖSUNG}}}
\rhead{\datum}

\lfoot{\zeit}
\cfoot{}
\rfoot{Seite \thepage\ von \pageref{LastPage}}

\newcommand{\punktebox}[1]{\hfill\fbox{\small \_\_\_/#1}}
\newcommand{\luecke}[1]{\rule{#1}{0.4pt}}
\newcommand{\lsg}[1]{\textcolor{loesung}{\textbf{#1}}}

\newtcolorbox{infobox}{
    colback=hellgrau,
    colframe=\fachfarbe,
    boxrule=1pt,
    arc=0pt,
    left=8pt, right=8pt, top=8pt, bottom=8pt
}

\newenvironment{aufgabe}[1]{%
    \needspace{5\baselineskip}%
    \nopagebreak[4]%
    \vspace{10pt}
    \noindent\textbf{\textcolor{\fachfarbe}{Aufgabe #1}}
    \nopagebreak[4]%
    \vspace{5pt}
    \nopagebreak[4]%
    \begin{list}{}{\leftmargin=0pt}
    \item[]
}{%
    \end{list}
}

\newcommand{\es}[1]{\textcolor{\fachfarbe}{\textbf{#1}}}

\begin{document}

% ═══════════════════════════════════════════════════════════════
% KOPFBEREICH
% ═══════════════════════════════════════════════════════════════

\begin{center}
    {\Large\bfseries\textcolor{\fachfarbe}{\thema}}\\[0.3em]
    {\large\textcolor{loesung}{\textbf{--- MUSTERLÖSUNG ---}}}
\end{center}

\vspace{5pt}

\noindent
\begin{tabularx}{\textwidth}{|X|X|X|}
\hline
\textbf{Name:} \lsg{Musterschüler/in} &
\textbf{Klasse:} \klasse &
\textbf{Datum:} \datum \\
\hline
\end{tabularx}

\vspace{10pt}

\begin{infobox}
\textbf{Zeit:} \zeit{} | \textbf{Punkte:} \maxpunkte{} BE | \textbf{Hilfsmittel:} keine
\end{infobox}

\vspace{15pt}

% ═══════════════════════════════════════════════════════════════
% AUFGABE 1: Orte in der Stadt
% ═══════════════════════════════════════════════════════════════

\begin{aufgabe}{1 -- Orte in der Stadt}
Übersetze die Orte ins Spanische (mit Artikel). \punktebox{6}
\end{aufgabe}

\noindent
\begin{tabular}{p{0.45\textwidth}p{0.45\textwidth}}
a) der Park = \lsg{el parque} & b) die Apotheke = \lsg{la farmacia} \\[0.8em]
c) der Bahnhof = \lsg{la estación} & d) das Kino = \lsg{el cine} \\[0.8em]
e) die Bibliothek = \lsg{la biblioteca} & f) der Supermarkt = \lsg{el supermercado} \\
\end{tabular}

\vspace{15pt}

% ═══════════════════════════════════════════════════════════════
% AUFGABE 2: ir + al/a la
% ═══════════════════════════════════════════════════════════════

\begin{aufgabe}{2 -- ir + al/a la}
Ergänze die richtige Form von \es{ir} und \es{al} oder \es{a la}. \punktebox{6}
\end{aufgabe}

\noindent
a) Yo \lsg{voy} \lsg{al} cine.

\vspace{0.8em}

\noindent
b) Tú \lsg{vas} \lsg{a la} biblioteca.

\vspace{0.8em}

\noindent
c) Nosotros \lsg{vamos} \lsg{al} parque.

\vspace{0.5em}
\textcolor{loesung}{\textit{Regel: a + el = al (Kontraktion)}}

\vspace{15pt}

% ═══════════════════════════════════════════════════════════════
% AUFGABE 3: estar + Ortsangaben
% ═══════════════════════════════════════════════════════════════

\begin{aufgabe}{3 -- estar + Ortsangaben}
Ergänze die richtige Form von \es{estar} und die Ortsangabe. \punktebox{6}
\end{aufgabe}

\noindent
a) El museo \lsg{está} \lsg{al lado de} la plaza. (neben)

\vspace{0.8em}

\noindent
b) La farmacia \lsg{está} \lsg{enfrente de} el banco. (gegenüber)

\vspace{0.8em}

\noindent
c) ¿Dónde \lsg{está} el hospital? -- \lsg{Cerca} del parque. (in der Nähe)

\vspace{15pt}

% ═══════════════════════════════════════════════════════════════
% AUFGABE 4: Wegbeschreibung
% ═══════════════════════════════════════════════════════════════

\begin{aufgabe}{4 -- Wegbeschreibung}
Ergänze die Richtungsangaben auf Spanisch. \punktebox{6}
\end{aufgabe}

\noindent
a) Geh geradeaus! = ¡Sigue \lsg{todo recto}! \hfill (2P)

\vspace{0.8em}

\noindent
b) Biege nach rechts ab! = ¡Gira \lsg{a la derecha}! \hfill (2P)

\vspace{0.8em}

\noindent
c) Biege nach links ab! = ¡Gira \lsg{a la izquierda}! \hfill (2P)

\vspace{15pt}

% ═══════════════════════════════════════════════════════════════
% AUFGABE 5: Dialog ergänzen
% ═══════════════════════════════════════════════════════════════

\begin{aufgabe}{5 -- Dialog ergänzen}
Ergänze den Dialog mit den richtigen Fragen. \punktebox{6}
\end{aufgabe}

\begin{tcolorbox}[colback=hellgrau, colframe=grau, boxrule=0.5pt, arc=0pt]
\noindent
\textbf{A:} Perdona, \lsg{¿dónde está la estación?} \hfill (2P)

\vspace{0.5em}

\noindent
\textbf{B:} La estación está cerca de aquí.

\vspace{0.8em}

\noindent
\textbf{A:} \lsg{¿Cómo llego a la estación?} \hfill (2P)

\vspace{0.5em}

\noindent
\textbf{B:} Sigue todo recto y gira a la derecha.

\vspace{0.8em}

\noindent
\textbf{A:} \lsg{¿Está lejos?} \hfill (2P)

\vspace{0.5em}

\noindent
\textbf{B:} No, está muy cerca. Cinco minutos.
\end{tcolorbox}

% ═══════════════════════════════════════════════════════════════
% PUNKTEÜBERSICHT
% ═══════════════════════════════════════════════════════════════

\vspace{25pt}
\hrule
\vspace{10pt}

\noindent
\begin{tabularx}{\textwidth}{|X|c|c|c|c|c||c|}
\hline
\textbf{Aufgabe} & 1 & 2 & 3 & 4 & 5 & \textbf{Gesamt} \\
\hline
\textbf{max. Punkte} & 6 & 6 & 6 & 6 & 6 & \textbf{30} \\
\hline
\textbf{erreicht} & \lsg{6} & \lsg{6} & \lsg{6} & \lsg{6} & \lsg{6} & \lsg{30} \\
\hline
\end{tabularx}

\vspace{10pt}

\begin{flushright}
\textbf{Note:} \lsg{14 NP}
\end{flushright}

\vspace{1em}

\textbf{Notenspiegel (14-NP-Skala):}

\begin{center}
\footnotesize
\begin{tabular}{|c|c|c|c|c|c|c|c|c|c|c|c|c|c|c|}
\hline
\textbf{NP} & 14 & 13 & 12 & 11 & 10 & 9 & 8 & 7 & 6 & 5 & 4 & 3 & 2 & 1 \\
\hline
\textbf{ab BE} & 30 & 29 & 27 & 26 & 24 & 22 & 20 & 18 & 16 & 14 & 12 & 10 & 8 & 6 \\
\hline
\end{tabular}
\end{center}

\vspace{1em}

\textcolor{loesung}{\textbf{Zusammenfassung Sequenz 6:}}
\begin{itemize}
    \item \textcolor{loesung}{\textit{ir} = gehen/fahren (Bewegung), \textit{estar} = sich befinden (Lage)}
    \item \textcolor{loesung}{\textit{a + el = al} (Kontraktion)}
    \item \textcolor{loesung}{\textit{¿Dónde está...?} (Wo?), \textit{¿Cómo llego a...?} (Wie komme ich zu?)}
\end{itemize}

\end{document}
